\section{The mean value therorem}
IF $x(t)$ is continuous on $a<=t<=b$ and differentiable on $a<t<b$ THEN
\[
	\frac{x(b)-x(a)}{b-a}
\]
for some $c$ with $a<c<b$ 

Geometric Interpretation 
\begin{figure}[H]
\centering
\def\svgwidth{\columnwidth}
\import{./images/}{mean_value_theorem.pdf_tex}
\end{figure}

\subsection{Least upper and Greatest lower Bounds} 
if $m$ is lower bound and $M$ is upper bound then,
\[
	m<=f(x)<=M
\]

\subsection{Bounding on average rate of change}
\[
	min_{a<=t<=b} \; x^{(1)}(t) <= \frac{x(b)-x(a)}{b-a} <= max_{a<=t<=b}\; x^{(1)}
\]

\section{Differential and Antiderivative}
consider a function $y=f(x)$ , Then the differential of y is 
defined as 
\[
	dy = F^{'}(x) \; d(x)
\]
This is also called the differential of $F$ and also denoted by $dF$. 
Rearranging this equation we get leibniz notation for Derivative
\[
	F^{'}(x) = \frac{dy}{dx} \; \text{or} \: \frac{dF}{dx}
\]
Explanation: 
$dy$ and $dx$ are known as infinitesimals . The point is that 
even though $dy$ and $dx$ are infinitely small quantities , Their ratio 
is NOT . The ratio is the derivative $F^{'}(x)$. $F^{'}(x)$ denotes 
How much $y$ changes if $x$ changes by tiny tiny bit. The geometric 
picture of differential is same as that of linear approximation. 

\begin{figure}[H]
  \centering
  \def\svgwidth{0.5\columnwidth}
  \import{./images/}{differential.pdf_tex}
\end{figure}

Linear approximation at $x$ $\Delta F \approx \dot F(x) \Delta x $
here $\Delta x$ is finite change in $x$ 

In terms of differential notation in $dF = \dot{F} dx$ 
here $dx$ denotes tiny tiny .. bit of $x$ i.e infinitesimals of $x$. 

\subsection{Antiderivative} 

\begin{definition*}
Antiderivative of $f(x)$ is any function $F(x)$ such that 
    \begin{equation*}
        \dot{F}(x) = f(x)
    \end{equation*}
\end{definition*}
Indefinite integral of $f(x)$ is denoted by $\int f(x)\: dx$. It is 
the family of functions. $\int f(x)\; dx = F(x) + C $  where $C$ is 
constant. 

Note : The constant of integration appears in this definition because 
derivative of function determines only the shape if function . But derivative does not change if the shape of function is shifted up or down by same
constant everywhere. 

\subsubsection{Uniqueness of Indefinite Integral}
The indefinite integral $\int f(x) \; dx = F(x) + C$ is termed indefinite
since it contains undetermined constant $C$ and it is the family of 
infinitely many function . But the constant is the only ambiguity of 
indefinite integral due to Mean value theorem Which guarranties that any
two Antiderivative of same function only differs by constant. 

\subsubsection{First rules of integration} 
\begin{tabular}{ll}
    \toprule 
    Differentiation rules & Integration rules \\
    \midrule
    $\int k f(x) = k \int f(x) dx$ & $d(kf) = k + f$ \\
    Sum \\
    $\int (f(x)+g(x)\; dx) = \int f(x) dx + \int g(x) dx$ & $d(F+G) =dF + dG$
\end{tabular}

The following naive product and quotient do not work . 
\begin{enumerate}

    \item $\int f * g$ DOES NOT EQUALL $\int f dx * \int g dx$
	    \item $\int \frac{f}{g}$ DOES NOT EQUALL $\frac{f dx}{g dx)}$

\end{enumerate}

\subsection{Method of Substitution} 
\paragraph{Hyperbolic Function}
\begin{enumerate}

  	\item 
	    \begin{equation*}
		    cosh(x) = \frac{e^{x}+e^{-x}}{2}
	    \end{equation*}
        \item 
	    \begin{equation*}
		sinh (x) = \frac{e^{x}-e^{-x}}{2}
	    \end{equation*}
	\item 
		\begin{equation*}
		    tanh(x) = \frac{sinh(x)}{cosh(x)}
		\end{equation*}
	\item 
		\begin{equation*}
		    cosh^{2}(x)- sinh^{2}(x) = 1
		\end{equation*}
\end{enumerate}

\section{ The Differential Equation} 
Differential equation involves differential coefficients. 
\begin{gather}
    \frac{dy}{dx} = f(x) \; \text{equation} \\
	y = \int f(x) \; dx \; \text{solution}
\end{gather}
 
\begin{tabular}{ll}
    \toprule
    equation     &    Differential equation \\
    \midrule
    $x^{2} -1=0$ & $x = -1 , +1$ $2$ points\\
    $\frac{dy}{dx}= f(x)$ & $ y  = \int f(x) dx $ family of functions \\
    \bottomrule
\end{tabular}

\subsection{Solving by seperation of variable} 
A differential equation  is seperable if it can be written in form . 
\begin{equation*}
    \frac{dy}{dx} = f(x) g(y)
\end{equation*}

can be solved as follows :  

\begin{gather}
    \frac{dy}{g(y)} = f(x) \; dx \\
    \int \frac{1}{g(y)} \; dy = \int f(x) \; dx
\end{gather}

\subsection{Slope fields}
Slope field is obtained as follows . At each point (x,y) you draw a short 
segment whose slope is the value of $\dot y$ at point $(x,y)$. The solution 
curves must be tangent to slope fields at all points. 

\paragraph{Euler's method} 
... see 18.03x
